% =====================================================================
% Abstract  (~180 words)
% =====================================================================
\begin{abstract}
This paper describes the design, implementation and evaluation of a
wireless 4 degree-of-freedom (DOF) robotic arm with gripper, developed
for the EE6003 Embedded Systems coursework at Dyson Institute of
Engineering + Technology. The system is built around two embedded
units: a handheld joystick controller based on an Arduino Uno R4
Minima, and a 3D-printed articulated arm driven by an Arduino Mega
2560. The two units communicate over a Bluetooth Classic SPP link
implemented with a paired HC-05 transmitter/receiver, transporting a
seven-byte custom-framed protocol with XOR checksum, start-byte
re-synchronisation and a link-loss watchdog. Each joint is closed
around a potentiometer feedback channel sampled by the Mega's ADC,
giving sub-degree position holding under payload disturbance. A
secondary on-board computer-vision subsystem --- initially an ESP32
WROVER CAM and later replaced by a Raspberry Pi camera module
communicating over UART --- enables a Sense--Think--Act autonomous
mode in which the arm scans, detects, classifies and sorts coloured
objects until at least five items have been placed. The report
explicitly separates core and advanced features in line with the brief
and presents results against accuracy, latency, repeatability and
robustness metrics.
\end{abstract}

\begin{IEEEkeywords}
embedded systems, robotic arm, Bluetooth, HC-05, closed-loop control,
finite state machine, computer vision, real-time, ISR, UART protocol
\end{IEEEkeywords}
